\chapter{Manual de usuario}
\label{anexo:manualusuario}

Este apéndice describe el uso de la herramienta desde el editor de Unity, orientado tanto a perfiles técnicos como a diseñadores de niveles y artistas.
Se recorren los componentes que el sistema pone a disposición del usuario, su configuración en el inspector y las ayudas visuales que facilitan la comprensión y la depuración de la navegación.
Para la instalación previa del paquete, consúltese el Apéndice~\ref{anexo:instalacion}.

\section{El componente NavVolume Space}

El componente \textit{NavVolume Space} es el punto de entrada del sistema: define la región del espacio sobre la que se calculará la navegación y centraliza su configuración.
Se añade a un objeto vacío de la escena desde el menú de componentes y expone, en su inspector, los parámetros fundamentales del volumen: su tamaño, el número de capas (que determina la resolución), la máscara de capas de física que se consideran obstáculos y el modo de construcción (véase la Figura~\ref{fig:manual-space}).

\pendientefig{Captura del inspector del componente \textit{NavVolume Space}, con las secciones de configuración (tamaño del volumen, número de capas y máscara de colisión), el modo de construcción y el panel de estadísticas desplegados.}{Inspector del componente \textit{NavVolume Space}.}{fig:manual-space}

\section{Definición del volumen en la escena}

El tamaño y la posición del volumen pueden ajustarse cómodamente desde la vista de escena.
Al activar la edición, aparecen unos manejadores (\textit{handles}) que permiten redimensionar la caja del volumen de forma interactiva, sin necesidad de introducir valores numéricos a mano (véase la Figura~\ref{fig:manual-handles}).

\pendientefig{Captura de la vista de escena mostrando el volumen de navegación delimitado por su caja envolvente y los manejadores interactivos para redimensionarlo.}{Definición y redimensionado del volumen en la vista de escena.}{fig:manual-handles}

\section{Construcción y precálculo (\textit{bake})}

El sistema ofrece tres modos de construcción del volumen, seleccionables en el inspector: \textit{precalculado} (\textit{Baked}), que carga una estructura previamente generada y almacenada en disco; \textit{construido al inicio} (\textit{BuildOnAwake}), que la genera al arrancar la escena; y \textit{manual}.
En el modo precalculado, un botón de \textit{Bake} lanza la construcción y muestra una barra de progreso cancelable mientras se procesa la escena (véase la Figura~\ref{fig:manual-bake}).

\pendientefig{Captura del botón de \textit{Bake} en el inspector y de la barra de progreso que se muestra durante la construcción del volumen.}{Proceso de construcción (\textit{bake}) del volumen.}{fig:manual-bake}

\section{Visualización del SVO mediante Gizmos}

Para facilitar la comprensión y la depuración del sistema, al seleccionar el volumen se representa su estructura interna en la vista de escena mediante Gizmos.
Un código de colores distingue los nodos vacíos, los que contienen geometría, las hojas parcialmente ocupadas y los \glspl{voxel} ocupados, lo que permite verificar de un vistazo que la ocupación calculada se corresponde con los obstáculos de la escena (véase la Figura~\ref{fig:manual-gizmos}).

\pendientefig{Captura de la visualización del \acrshort{svo} en la vista de escena mediante Gizmos, sobre un escenario con obstáculos, mostrando el código de colores que distingue los distintos tipos de nodo y los vóxeles ocupados.}{Visualización del \acrshort{svo} mediante Gizmos.}{fig:manual-gizmos}

\section{Panel de estadísticas}

El inspector incorpora un panel de estadísticas que ayuda a dimensionar y ajustar el volumen.
En él se muestran el tamaño de vóxel resultante, el número de \glspl{voxel} reservados frente al total teórico (con su porcentaje de ahorro), una estimación de la memoria ocupada y un desglose de los tiempos de la última construcción (véase la Figura~\ref{fig:manual-stats}).

\pendientefig{Captura del panel de estadísticas del inspector: tamaño de vóxel, vóxeles reservados frente al total teórico, memoria estimada, gráfico de nodos por capa y desglose de tiempos de construcción.}{Panel de estadísticas del volumen.}{fig:manual-stats}

\section{El componente NavVolume Agent}

El componente \textit{NavVolume Agent} dota a un objeto de la capacidad de navegar de forma autónoma por el volumen.
Su inspector permite configurar los parámetros de movimiento del agente y asociarlo a un tipo de agente concreto (véase la Figura~\ref{fig:manual-agent}).
Durante la ejecución, la trayectoria calculada para el agente puede representarse mediante Gizmos, lo que facilita comprobar visualmente el resultado de la búsqueda y del suavizado (véase la Figura~\ref{fig:manual-agent-path}).

\pendientefig{Captura del inspector del componente \textit{NavVolume Agent} con sus parámetros de movimiento y de configuración.}{Inspector del componente \textit{NavVolume Agent}.}{fig:manual-agent}

\pendientefig{Captura de la vista de escena mostrando, mediante Gizmos, la trayectoria calculada para un agente desde su posición hasta su destino.}{Visualización de la trayectoria de un agente mediante Gizmos.}{fig:manual-agent-path}

\section{El componente NavVolume Obstacle}

El componente \textit{NavVolume Obstacle} marca a un objeto como obstáculo dinámico, de modo que los agentes lo tengan en cuenta durante la evasión local en tiempo de ejecución.
Su inspector permite ajustar la forma y las dimensiones del obstáculo (véase la Figura~\ref{fig:manual-obstacle}).

\pendientefig{Captura del inspector del componente \textit{NavVolume Obstacle}, con la configuración de su forma y dimensiones.}{Inspector del componente \textit{NavVolume Obstacle}.}{fig:manual-obstacle}

\section{Tipos de agente}

Los tipos de agente se definen como recursos reutilizables (\textit{assets}) que encapsulan los parámetros comunes a un conjunto de agentes, como su radio.
De este modo, un mismo volumen puede atender a varias clases de agente, y cada agente se asocia al tipo que le corresponde (véase la Figura~\ref{fig:manual-agenttype}).

\pendientefig{Captura del inspector de un recurso de tipo de agente (\textit{Agent Type}) mostrando sus parámetros configurables.}{Inspector de un tipo de agente.}{fig:manual-agenttype}

\section{Uso desde código}

Más allá de la configuración visual, el sistema expone una interfaz de programación para los usuarios que requieran un control más fino.
A través de ella es posible solicitar rutas, comprobar si un punto del espacio es navegable, ajustar una posición al punto navegable más cercano u obtener puntos navegables aleatorios, entre otras operaciones.
La documentación detallada de esta interfaz queda fuera del alcance de este manual.
